% Proposed B.136. Isolated owned fragment; central integration remains root-owned.
\subsection*{The actual unfilled edge groups and their symmetric local filling}

\statement{B.136}{Proved}{The unfilled inclusion edge embeds, its target contains a full free-six subgroup, and the proper local filling has an explicit fixed-point-free involution}
\label{stmt:B136}
Use the actual B.134 marking
\[
 V=R^0*_{K}p(R^0),\quad K=F(a,B,d,e,f),\quad K_1=CKC,
 \quad L=\langle K,k\rangle,\quad E_1=CLC,
\]
where $p|_K=\eta$ and $\theta$ interchanges the two vertices.
Put $h=fk=Cz$, $k_1=CkC=hkh^{-1}$, $f_1=CfC$ and
\begin{equation}\label{eq:B136-words}
 \begin{gathered}
 a_1=CaC,\qquad w_0=aa_1^{-1},\qquad v_0=p(w_0),\\
 s'=w_0^{-1}p(k)w_0,\qquad s=v_0^{-1}kv_0,\qquad
 \sigma(CuC)=p(C\eta(u)C)\quad(u\in K).
 \end{gathered}
\end{equation}
Then the actual regional intersections are
\begin{equation}\label{eq:B136-regional-intersections}
 \begin{gathered}
 E_1\cap K=\langle f\rangle,\qquad K_1\cap K=1,\\
 \widehat E_1\cap\widehat K=\overline{\langle f\rangle},
 \qquad \widehat K_1\cap\widehat K=1,
 \end{gathered}
\end{equation}
where the completed groups are their actual embedded closures in
$\widehat{R^0}$.

The inclusion-side map of the abstract B.134 edge
\[
 i_0:E=E_1*\langle y\rangle\longrightarrow V,
 \qquad i_0|_{E_1}=\mathrm{id},\quad i_0(y)=s'
\]
is injective. Its mirror is the embedded group
\begin{equation}\label{eq:B136-two-edges}
 \theta i_0(E)=\langle p(E_1),s\rangle=p(E_1)*\langle s\rangle.
\end{equation}
The actual subgroup
\begin{equation}\label{eq:B136-target-free}
 T_6=\langle\sigma(K_1),s\rangle
       =\sigma(K_1)*\langle s\rangle\cong F_6
\end{equation}
has its full induced profinite topology in $V$.
No full-topology assertion for $i_0(E)$ is inferred merely from
its ordinary embedding.

Write $\eta_1=\operatorname{Ad}(C)\eta\operatorname{Ad}(C)$ on $K_1$,
and put
\[
 (u_1,u_2,u_3,u_4)
    =\bigl(CaC,CdC,CeC,C(f^{-1}Bf)C\bigr).
\]
In $F_*=K_1*\langle k_1\rangle*\langle y\rangle$ let
\begin{equation}\label{eq:B136-local-relators}
 \mathcal R=
 \bigl\{[(f_1k_1)^2,k_1],\ [k_1,u_i]\ (1\leq i\leq4)\bigr\}.
\end{equation}
The automorphism $\kappa_0$ of $F_*$ acts by $\eta_1$ on $K_1$
and exchanges $k_1,y$. It is an involution. The group
\begin{equation}\label{eq:B136-Dloc}
 D_{\mathrm{loc}}
 =F_*/\langle\!\langle\mathcal R,\kappa_0(\mathcal R)
                                      \rangle\!\rangle
\end{equation}
is a proper ordinary amalgam of two copies of $E_1$ over $K_1$,
with the second edge marking twisted by $\eta_1$. It has a proper
full profinite amalgam completion and full vertex topologies.
The induced involution $\kappa$ satisfies
\begin{equation}\label{eq:B136-fixed}
 D_{\mathrm{loc}}^{\kappa}=1,
 \qquad (\widehat D_{\mathrm{loc}})^{\kappa}=1.
\end{equation}

Let $V_*=V/N_*$ be the actual B.134 quotient by all five
$\rho_i$ and their five $\theta$-mates. There is an actual map
$a_{\mathrm{loc}}:D_{\mathrm{loc}}\to V_*$, and the other edge map is
\begin{equation}\label{eq:B136-two-maps}
 b_{\mathrm{loc}}=\theta a_{\mathrm{loc}}\kappa.
\end{equation}
Replacing $E$ by $D_{\mathrm{loc}}$ in B.134's pushout does not
change the actual group. Its initial ordinary kernel-matching
condition is exactly
\begin{equation}\label{eq:B136-kernel-test}
 \ker a_{\mathrm{loc}}\text{ is }\kappa\text{-invariant}.
\end{equation}
The matching finite-index normal subgroups of $V_*$ may be
chosen $\theta$-invariant cofinally. For such a subgroup $N$,
the finite matching condition is exactly
\begin{equation}\label{eq:B136-finite-test}
 a_{\mathrm{loc}}^{-1}(N)\text{ is }\kappa\text{-invariant}.
\end{equation}
Neither \eqref{eq:B136-kernel-test} nor an embedding of
$D_{\mathrm{loc}}$ in $V_*$ is asserted. The fixed-subgroup
statement \eqref{eq:B136-fixed} is strictly local.

\begin{proof}
\emph{B.136.1. Retractions and all-modulus lamp detection.}
The actual ordinary retraction $r:R^0\to K$ of B.112/B.117 fixes
$K$ and satisfies
\[
 r(CuC)=\zeta(u),\quad
 \zeta(a)=\zeta(B)=\zeta(d)=\zeta(e)=1,\quad\zeta(f)=f,
 \qquad r(k)=r(k_1)=1.
\]
Thus $r(E_1)=\langle f\rangle$. Moreover $f$ belongs to $E_1$:
$ChC=h$, so $E_1$ contains $h,k_1$ and hence
$k=h^{-1}k_1h$ and $f=hk^{-1}$. A common element of $E_1$ and
$K$ is fixed by $r$, and therefore lies in $\langle f\rangle$.
This proves the first intersection. Both subgroups have full
induced topology by B.117 and conjugation. The identical argument
with the continuous retraction proves its completed counterpart;
no completion of an intersection is substituted.

For the second intersection, a common element of $\widehat K$
and $\widehat K_1$ is $f^{\lambda}$, by $r$. Applying the
$C$-conjugate retraction shows that it is also $(CfC)^{\mu}$.
The actual $z$-height character in the full B.117 regional model
\[
 A(\Omega)\rtimes
       \bigl((C_2(C)*\mathbb Z(k))\times\mathbb Z(z)\bigr)
\]
sends both $f=Ck^{-1}z$ and $CfC=k^{-1}Cz$ to $1$.
Consequently $\lambda=\mu$ in $\widehat{\mathbb Z}$.

For every integer $n\geq2$, kill all graph generators and $z$,
and map $C$ to the lamp at $0$ and $k$ to the shift in
$C_2\wr C_n$. This is an actual finite regional homomorphism.
The two words have images $(\epsilon_0,-1)$ and
$(\epsilon_{-1},-1)$. For any integer $L$, the lamp vectors of
their $L$th powers differ by
\[
 \epsilon_0+\epsilon_{-L}.
\]
This follows by telescoping successive shifted lamps, for positive
and negative $L$. Both elements have order dividing $2n$, so the
same formula applies to a profinite exponent by choosing its
representative modulo $2n$. Equality forces $L=0\pmod n$.
Thus $\lambda$ is zero modulo every $n$, hence is zero. This
proves the second completed intersection, and also the ordinary
one in \eqref{eq:B136-regional-intersections}.

\emph{B.136.2. The actual free-six target and its full topology.}
Conjugating $T_6$ by $v_0$ gives
\begin{equation}\label{eq:B136-target-conjugation}
 v_0T_6v_0^{-1}=\langle D',k\rangle,
 \qquad D'=p(aK_1a^{-1}).
\end{equation}
Indeed $w_0C=aCa^{-1}$ and $\eta(K)=K$, so
\[
 v_0\sigma(K_1)v_0^{-1}
     =p(aCa^{-1}KaCa^{-1})=p(aK_1a^{-1}).
\]
The right-vertex subgroup $D'$ has trivial intersection with $K$
by B.136.1; $a\in K$ and $p$ preserves $K$. The left-vertex
subgroup $\langle k\rangle$ has trivial intersection with $K$
by actual base winding. Ordinary amalgam normal form proves
\eqref{eq:B136-target-free}.

We spell out the requisite completion of $V$. Given finite
quotients of its two $R^0$ vertices, choose a finite-index fully
invariant subgroup of the identified $K$ contained in their
edge restrictions, using the actual $\eta$ identification.
For each vertex combine its prescribed finite quotient with
its actual retraction onto that identified $K$, followed by
this common quotient. The exact edge kernels are then equal.
The resulting finite vertex groups admit faithful regular
permutation actions on one common finite set, after taking
suitable multiples. Their restrictions to the finite edge group
are free with equal orbit counts, so identify the orbits
equivariantly. This gives a finite quotient of the amalgam
retaining both vertex quotients. It proves cofinal full vertex
and edge topology. The universal property identifies
\[
 \widehat V=\widehat{R^0}
       \amalg_{\widehat K}p(\widehat{R^0})
\]
as a proper full profinite amalgam.

Both factors in \eqref{eq:B136-target-conjugation} have their
full topologies in their vertices. Their actual closed
intersections with $\widehat K$ are trivial, by the completed
intersection in B.136.1 and the continuous base-winding character.
The closed-vertex-subgroup theorem
\cite[Corollary~2.18, with Lemma~2.17]{BLopesZalesskii2025}
therefore identifies their closed generated subgroup with their
free profinite product. It is $\widehat F_5\amalg
\widehat{\mathbb Z}=\widehat F_6$. Conjugating back proves the
full topology of $T_6$. The zero-edge-intersection hypotheses
have been verified, rather than replaced by separate vertex
embeddings alone.

\emph{B.136.3. The inclusion edge and its actual branch normal form.}
Conjugation by $w_0$ changes $E_1$ to $aE_1a^{-1}$, since
$a_1\in E_1$, and changes $s'$ to $p(k)$. Put
$\ell=afa^{-1}$. By B.136.1 the left subgroup has edge
intersection exactly $\langle\ell\rangle$. On the right,
$\langle\ell,p(k)\rangle$ is the image under $p$ of
$\langle\eta(\ell),k\rangle$.

Use the complete marked B.117 amalgam
\begin{equation}\label{eq:B136-L-amalgam}
 \begin{gathered}
 L=H_{\mathrm{ret}}*_{\langle k\rangle}P_{\mathrm{leaf}},\qquad
 H_{\mathrm{ret}}=\langle h,k\mid[h^2,k]\rangle,\\
 P_{\mathrm{leaf}}=\langle k\rangle\times F(a,d,e,V_4),
 \qquad f=hk^{-1},\quad V_4=f^{-1}Bf.
 \end{gathered}
\end{equation}
The certified literal reflection gives
\begin{equation}\label{eq:B136-branch-word}
 \begin{gathered}
 \eta(\ell)=wf^{-1}w^{-1},\\
 w=afV_4f^{-1}a^{-1}f(de d^{-1})f^{-1}d^{-1}.
 \end{gathered}
\end{equation}
Each displayed free-leaf block is nontrivial. Every nonzero
$f$-power is outside $\langle k\rangle$, by the $h$-exponent
character of $H_{\mathrm{ret}}$. For every nonzero integer $n$,
also $f^{-1}k^nf\notin\langle k\rangle$: the quotient
$H_{\mathrm{ret}}\to C_2(C)*\mathbb Z(k)$, $h\mapsto C$,
sends it to the reduced word $kCk^nCk^{-1}$ outside that cyclic
subgroup.

Expand $w^{-1}k^nw$ in \eqref{eq:B136-L-amalgam}.
The first leaf $a$ commutes with $k$, leaving central block
$f^{-1}k^nf$. All its adjacent leaf blocks and the remaining
alternating $f$ and leaf blocks are nontrivial outside the edge.
The reduced word begins in the leaf factor with $d$ and ends
there with $d^{-1}$. Consequently, for all nonzero exponents,
\[
 (w^{-1}k^{n_1}w)f^{m_1}\cdots
 (w^{-1}k^{n_r}w)f^{m_r}
\]
is reduced and nontrivial. The joining $f^{m_i}$ blocks are
outside the edge, between the nontrivial leaf end blocks.
Pure powers are nontrivial as well. This proves
\[
 \langle k,\eta(\ell)\rangle
     =\langle k\rangle*\langle\eta(\ell)\rangle.
\]
The actual retraction $L\to K$ kills $k$ and fixes
$\eta(\ell)$. Hence an element of this subgroup that lies in
$K$ equals its retraction and is a power of $\eta(\ell)$.
Applying $p$ gives
\[
 \langle\ell,p(k)\rangle
      =\langle\ell\rangle*\langle p(k)\rangle,
 \qquad \langle\ell,p(k)\rangle\cap K=\langle\ell\rangle.
\]
Thus both vertex subgroups have exactly the same intersection
with the ambient edge. Ordinary reduced-word normal form
embeds their subamalgam, and gives
\[
 \langle aE_1a^{-1},p(k)\rangle
 =aE_1a^{-1}*_{\langle\ell\rangle}
      (\langle\ell\rangle*\langle p(k)\rangle)
 =aE_1a^{-1}*\langle p(k)\rangle.
\]
Conjugating by $w_0^{-1}$ proves injectivity of $i_0$.
Applying $\theta$ proves \eqref{eq:B136-two-edges}.
This argument does not assert a profinite subamalgam theorem
for the nontrivial cyclic intersection.

\emph{B.136.4. Both local fillings and their exact involution.}
B.117 gives $E=F_*/\langle\!\langle\mathcal R\rangle\!\rangle$.
The B.134 words satisfy, literally in $i_0(E)$,
\[
 i_0\bigl(\kappa_0(\mathcal R)\bigr)
       =\{\theta(\rho_0),\ldots,\theta(\rho_4)\}.
\]
Indeed $\kappa_0$ applies $\eta_1$ to the five free edge
letters and exchanges $k_1,y$; the target substitutions are
exactly $\sigma$ and $s$, followed by $\theta$.
Since $\sigma(K_1)=p(K_1)$, the entire target generator group
is $\langle p(E_1),s\rangle$, including $p(k_1)$.

The first five relations in \eqref{eq:B136-Dloc} present one
copy of $E_1$. Applying the invertible $\eta_1$ basis change,
the other five present the second copy with base letter $y$.
This gives the stated amalgam over $K_1$. Each edge is an
actual retract, by B.117. Ordinary normal form therefore makes
the amalgam proper. For full completion, use a common fully
invariant finite $K_1$ kernel below two prescribed vertex
restrictions, combine each vertex target with its edge
retraction, and apply the regular-action construction in
B.136.2. This proves proper full completion and full vertex
topology. The involution $\kappa_0$ exchanges the two entire
relator sets, so induces the stated $\kappa$.

Both sets are killed by the actual $N_*$. Thus $i_0$ induces
$a_{\mathrm{loc}}$; evaluation on $K_1,k_1,y$ gives
\eqref{eq:B136-two-maps}. On these generators the right side
is respectively $\sigma$, $s$, and $p(k_1)$.
The two maps already kill both local relator sets. Factoring
through $D_{\mathrm{loc}}$ therefore leaves every original
edge equation and both cross equations unchanged. No
embedding of this local quotient in $V_*$ follows.

\emph{B.136.5. Exact ordinary and finite matching tests.}
Since $\theta$ is an automorphism,
\[
 \ker b_{\mathrm{loc}}=\kappa(\ker a_{\mathrm{loc}}).
\]
This proves \eqref{eq:B136-kernel-test}. When it holds, B.134's
matching union for these maps begins and remains trivial,
so $V_*$ embeds as the proper ordinary HNN base. Conversely,
a base embedding in the actual pushout forces equality of
the two edge kernels, hence forces the stated invariance.
No such invariance is proved here.

For a finite-index normal $N\triangleleft V_*$ write
$A_N=a_{\mathrm{loc}}^{-1}(N)$. Equation
\eqref{eq:B136-two-maps} gives
\[
 b_{\mathrm{loc}}^{-1}(N)=\kappa(A_{\theta(N)}).
\]
Thus matching is $A_N=\kappa(A_{\theta(N)})$, a condition
symmetric in $N,\theta(N)$. A matching $N$ can be refined by
$N\cap\theta(N)$, which is again matching and is
$\theta$-invariant. For these cofinal normals the condition
is exactly \eqref{eq:B136-finite-test}. This is B.134's
independent finite criterion with the actual local symmetry
made explicit; it does not equate a closed ordinary matching
union with the compact kernel.

\emph{B.136.6. The folded local group and its fixed subgroup.}
There are exact inverse presentations
\begin{equation}\label{eq:B136-fold}
 F_{\mathrm{fold}}=D_{\mathrm{loc}}\rtimes\langle\kappa\rangle
    =E_1*_{K_1}Q_1,
 \qquad Q_1=K_1\rtimes_{\eta_1}C_2(\kappa).
\end{equation}
Eliminate the second copy of $E_1$ by writing it as
$\kappa E_1\kappa$; the common-edge relation is precisely
$\kappa u\kappa=\eta_1(u)$. Conversely, the parity kernel
of the right-hand group has the two $E_1$ vertices with this
edge marking, recovering \eqref{eq:B136-Dloc}.

Here too full completion is proper, with a concrete finite
proof. Given finite targets of $E_1,Q_1$, choose an
$\eta_1$-invariant finite-index normal subgroup $L_0$ of $K_1$
inside both edge restriction kernels. Intersect the first
restriction with its $\eta_1$-image and with the second
restriction; the second is already invariant. Combine the
$E_1$ target with its actual retraction to $K_1/L_0$, and
combine the $Q_1$ target with $(K_1/L_0)\rtimes C_2$.
Both exact edge kernels are $L_0$. The faithful finite
regular-action amalgamation proves cofinal full topology
as before. Therefore
\[
 \widehat F_{\mathrm{fold}}
      =\widehat E_1\amalg_{\widehat K_1}\widehat Q_1
\]
is the proper full completion. Its edges are nondegenerate:
$k_1$ has nonzero base winding in $E_1$, while $\kappa$ is
outside the index-two $K_1$ in $Q_1$.

The certified actual prefix basis $P_1,\ldots,P_5$ of $K$
satisfies $\eta(P_i)=P_i^{-1}$. Its $C$-conjugate is a free
basis $q_1,\ldots,q_5$ of $K_1$ inverted by $\eta_1$.
Consequently $Q_1$ has the inverse presentation
\[
 Q_1=\langle\kappa,q_1,\ldots,q_5\mid
       \kappa^2=1,\ (\kappa q_i)^2=1\rangle
      \cong\mathop{*}_{i=0}^{5}C_2.
\]
The same presentation and full-completion universal property
give the free profinite product of six $C_2$ groups. Its
selected involution centralizer is exactly
$\langle\kappa\rangle$, by ordinary normal form and by
\cite[Theorem~B$'$]{A_HR1985}, respectively.

In the ordinary tree of \eqref{eq:B136-fold}, $\kappa$ fixes
the $Q_1$ vertex but no edge, since each edge stabilizer is a
conjugate of the torsion-free free group $K_1$. Its connected
fixed subtree is that single vertex. Every centralizer element
preserves it, so the whole centralizer is
$C_{Q_1}(\kappa)=\langle\kappa\rangle$.

The compact argument uses the actual completed tree. The free
profinite edge group $\widehat K_1$ has no involutions. Hence
$\kappa$ fixes no edge, while it fixes the displayed
$\widehat Q_1$ vertex. Its fixed subgraph is connected by
\cite[Theorem~2.8]{B_ZalesskiiMelnikov1989}. A connected
profinite graph without edges has just one vertex. Thus every
compact centralizer element belongs to $\widehat Q_1$, whose
centralizer is again $\langle\kappa\rangle$.
Intersecting these centralizers with the respective parity
kernels $D_{\mathrm{loc}}$ and $\widehat D_{\mathrm{loc}}$
proves \eqref{eq:B136-fixed}. The latter is the actual parity
kernel because the index-two subgroup has full topology.
No torsion-freeness of $\widehat E_1$ was assumed or needed.

All these assertions precede the ambient kernel of
$a_{\mathrm{loc}}$. Proper local amalgams and their fixed
subgroups do not prove their survival after both ambient
normal closures are imposed. An invariant quotient kernel can
create new fixed elements, as the separate actual inner-star
face demonstrates. The actual kernel symmetry, source
intersection, and selected fixed-source equality remain open.
\end{proof}

\begingroup\raggedright\footnotesize
The complete literal branch and basis certificate is
\path{work_DE/next_marginal_source/check_unfilled_edge_maps27.py}
with its genuine output and JSON. The all-power normal forms,
all-modulus lamp proof, finite-target constructions and exact
local symmetry are proved above and in the preserved records
\path{work_DE/next_marginal_source/unfilled_edge_map_structure27.md}
and
\path{work_DE/next_marginal_source/local_edge_involution_fixed_subgroup27.md}.
The two actual cross sheets of B.134 are retained throughout.
\par\endgroup
